\documentclass{article}
\usepackage{graphicx}
\usepackage{ctex}
\title{宽带高增益RFID标签天线的研究与设计}
\author{Katerina Svandova;Zdenek Smutny}
\date{2024年7月18日}

\begin{document}
	
	\maketitle
	
	\begin{abstract}
	射频识别(Radio Frequency Identification,RFID)作为物联网产业中的关键技术之一，是一种通过电磁波以非接触方式对目标信息进行自动识别的技术，其在诸多领域都具有广阔的应用前景。RFID标签天线是RFID系统的一个基本组成单元，因而标签的性能对RFID系统的整体性能有着显著的影响。为此，围绕实际工程应用需求，本文开展了宽带高增益RFID标签天线的研究与设计工作，主要的内容如下：
	1.基于开环谐振器(SplitRing Resonator,SRR)的RFID天线设计。根据SRR结构谐振的产生机理，分析了SRR内外环的半径和开口宽度等结构参数对天线性能的影响，并通过SRR内外环之间的相互耦合，在谐振点处实现了与偶极子天线类似的全向辐射特性。进而，设计了一款基于SRR的小型化RFID标签天线，仿真结果表明，该天线在中心谐振频率915MHz处具有110MHz的带宽，增益可达-1.3dB,且所具有的全向辐射特性能较好地满足诸多实际应用需求，但相较于偶极子天线，其增益并不占优势。
	2.弯折型偶极子RFID标签天线的设计。以偶极子作为基础结构，采用弯折线技术来实现标签天线的小型化，分析了对于谐振频率和输入阻抗影响较大的结构参数。而后，通过分析和优化不同的耦合结构，实现了标签天线与标签芯片之间的良好阻抗匹配，并以此为基础设计了多款加载耦合结构的弯折型偶极子RFID标签天线。最后，选择综合性能最优的设计进行实物制作，该款弯折型偶极子标签天线中心频率同样为915MHz,带宽为90MHz,天线的增益可达2.3dB,尺寸仅为25×52×1.6mm³,辐射方向图具有良好的全向性，且对比文献中已报道的偶极子RFID天线，其在增益及尺寸方面都更有优势。由于加工误差的存在，最后实测得到的工作频点为910MHz,该频点处标签的最大增益为1.8dB,同时在全方向上都有着较好的阅读距离，最大阅读距离为11.5m。
	3.基于机器学习的RFID标签天线设计方法研究。第一部分是正向设计，首先以对标签天线性能影响最大的结构参数(即天线末端展宽的宽度W和高度H)作为输入，天线的带宽作为输出，建立神经网络正向预测模型；然后，外接遗传优化算法(Genetic Algorithm,GA),将天线的带宽作为优化的目标函数，寻找其最大值；所得结果显示：通过基于机器学习模型的正向优化，可实现124MHz的工作带宽，并且与利用全波仿真软件的优化方法相比极大地节省了时间。第二部分是逆向设计，与正向设计方法不同，逆向设计是将期望的天线的带宽作为神经网络的输入，对应的天线结构参数W和H作为输出；这里通过将正向设计所得到的带宽作为期望输入，利用训练好的逆向设计网络给出对应天线的结构参数；仿真验证显示：逆向设计网络能够较为精确地预测出期望的天线结构，这初步验证了逆向设计的可行性。
	\end{abstract}
	
	\section{研究背景}
	射频识别(RFID)技术是一种无线识别技术，其基本原理是通过无线电频率，使阅读器与标签之间实现目标识别并获取其数据，被认为是21世纪最具发展潜力的信息技术之一。RFID技术在物流管理、交通、防伪以及身份识别等各个领域都有着广泛的应用。相较于二维码、条形码等其他的识别技术，RFID技术具有显著的优势，它不需要与物体直接接触也可以实现完全的自动识别[,可以适用于复杂的环境，拥有更远的读取距离和更高的读取效率。
	随着物联网和智能化技术的兴起，对于远程、快速、准确的物体识别与管理的需求日益增长。作为物联网的子行业，RFID技术被视为物联网发展的基础和前提。RFID技术与互联网技术的相互结合推动了RFID技术在不同领域的发展，随着智能化的不断发展RFID技术将在更广泛的领域发挥作用，为社会和经济发展带来更多机遇。未来，随着相关技术的完善和成熟，RFID技术在生产、生活领域中的重要性将进一步凸显[2]。
	常见的RFID系统主要由标签、阅读器以及对采集的数据信息进行分析处理的数据管理系统组成，相对于阅读器与数据管理系统，RFID标签是最为重要的部分，标签性能的优劣直接决定RFID系统的准确性，考虑到不同的应用场景，要求标签具备很强的抗干扰性以及稳定性，可以在复杂的环境下工作[3。
	本文主要针对UHF频段的RFID标签天线的设计，通过对天线结构的优化，提高天线的带宽和增益等性能。
	
	\section{研究方法}
    本文以RFID标签天线的设计为主要研究内容，结合天线的基础理论，针对实际的工程应用，着重于标签天线的设计与优化。本文主要分为以下六个章节进行叙述：
    第一章从国内外研究现状出发，首先介绍了RFID技术的发展背景，介绍了近年来国内外有关天线的小型化、高增益等优化设计的发展现状。并说明了机器学习在天线优化设计中的优势及发展。
    第二章首先对RFID的整体系统做一个概述，介绍了无源RFID标签天线的类型和特点；对天线的电磁学基础，以及天线的辐射场特性做了简单的阐述；介绍了描述标签天线性能的常用参数指标：阅读距离、带宽等。给出了天线在不同应用场景、需求下标签天线的设计特点及设计流程。
    第三章简要介绍了开环谐振器(SRR)的基础理论，研究了SRR的电磁特性，分析了SRR内外环的结构参数对天线性能的影响，验证了SRR在内外环的相互耦合下产生的第二谐振点所呈现的辐射特性与偶极子天线相似，并设计了一款基于SRR的小型化RFID标签天线，可以作为偶极子结构的替代结构用于实现天线结构的简单化、小型化。
    第四章以半波偶极子为基础，简述了偶极子天线的相关理论，利用电磁仿真软件HFSS分析了不同的结构参数对于偶极子天线的性能影响，确定了以弯折线结构为基础的天线结构。然后通过分析和优化不同的耦合结构，设计了多款加载耦合结构的弯折型偶极子RFID标签天线，将综合性能最优的一款天线制作成实物并测试，实测结果显示：天线具备良好的性能，完成了弯折型偶极子RFID标签天线的设计。
    第五章主要研究了与机器学习算法相结合的天线优化设计，以宽带宽为优化的目标，通过HFSS对偶极子天线参数进行仿真得到大量的仿真数据结果，并且结合遗传算法进行优化，得到天线最优的带宽值。再将机器学习预测的最优的带宽值作为输入，通过机器学习预测的模型，得到所需的天线结构参数，在HFSS中对得到的数据进行仿真验证，实现机器学习算法对天线的优化设计。
	\section{研究结果}
随着射频识别技术的不断发展，它广泛地应用于物流管理、交通等领域，作为 RFID
技术的基本组成单元，RFID 标签性能的优劣对整个 RFID 系统至关重要。本文主要的研
究内容是对高增益、宽带宽的 RFID 标签天线的研究与设计，设计并制作了一款应用于超
高频的标签天线，完成了标签指标测试；最后结合遗传算法，采用机器学习的方法来设计
和优化天线。
首先是从具有负磁导率、负介电常数的左手材料出发，提出了具有左手材料性质的
SRR 结构，通过其产生谐振的机理，分析了 SRR 的内外环半径及开口宽度等结构对天线
性能的影响，并通过内外环之间的相互耦合，激发出第二个谐振点，在该谐振点处实现了
与偶极子类似的全向辐射特性。进而，通过改变 SRR 的结构，设计了一款 RFID 标签天
线用于替代偶极子天线。该天线工作在 915MHz，带宽达到 110MHz，增益为-1.3dB，具
有的全向辐射性能较好的满足实际需求，但是相较于偶极子天线增益较低。
其次以偶极子作为标签天线的基础结构，介绍了偶极子天线上的电流分布及辐射特
性。采用弯折线技术实现偶极子天线的小型化，并分析了对于谐振频率和输入阻抗影响较
大的结构参数。由于标签天线需要与标签芯片直接连接，不能使用匹配电路，因此通过分
析和优化不同的耦合结构，实现了天线与芯片之间的良好阻抗匹配，在此基础上设计了多
款加载耦合结构的标签天线。最后，选择综合性能最优的加载短截线结构的弯折型偶极子
RFID 标签天线进行实物制作，该款天线中心谐振频率为 915MHz，带宽为 90MHz，最大
增益 2.3dB，尺寸仅为 25×52×1.6mm3，实际测得的工作频点为 910MHz，该频点处的最
大增益为 1.8dB，在全方向上都有较好的阅读距离，最大阅读距离为 11.5m。
最后将天线的优化设计与机器学习方法相结合。首先是正向的设计方法，将天线末端
展宽的宽度 W 和高度 H 作为神经网络模型的输入，天线的带宽作为输出，建立神经网络
正向预测模型；然后外接遗传优化算法，将天线的带宽作为优化的目标函数，寻找其目标
范围内的最大值，并将最大值所对应的 W 和 H 带入 HFSS 中仿真验证，所得结果显示：
基于机器学习模型的正向优化可以实现 124MHz 的工作带宽，与全波仿真软件的优化方
法对比，极大地节省了时间，优化效率提高了大约 130 倍。然后是逆向设计，与正向设计
不同，逆向设计是将期望的天线带宽作为神经网络输入，对应的天线结构参数 W 和 H 作
为输出，，建立逆向神经网络预测模型；将正向优化所得的带宽作为输入，利用训练好的
逆向预测模型给出对应的结构参数；仿真验证显示：逆向设计模型能够较为准确地预测出
期望的天线结构，验证了逆向设计的可行性
	
	\section{结论}
	本文主要开展了宽带高增益 RFID 标签天线的研究与设计，并且使用了机器学习的方
	法来优化天线，但是整个过程仍然存在一些问题需要解决和完善。首先是本文所设计的标
	签天线所处的环境是常见的微波暗室，没有考虑到复杂环境下的变化，比如在天线旁边放
	置金属会产生什么影响，后续可以在标签天线的抗干扰方向上继续深入研究。然后是机器
	学习优化天线方面，神经网络模型的训练集训练的样本数据较少，更多的数据会使神经网
	络模型更加的准确；且本文只针对偶极子天线这一种结构进行优化，后续可以尝试用机器
	学习方法优化更多不同类型的天线。
	
\end{document}